\section{The auxiliary polynomials and their determinant}
\label{sec:algebra}

\subsection{Intervals, poles, and normalization}

Fix the parameters of Theorem~\ref{thm:main}. Let \(n\) run through
positive multiples of \(L\). List all \(M=m+1\) intervals, including
repetitions, in nested order:
\[
I_a=[\alpha_a,\beta_a],\quad 0\le a\le m,\qquad
0=\alpha_0\le\cdots\le\alpha_m<n/2,\quad
\beta_a=n-\alpha_a.
\]
Thus \(I_m\subseteq\cdots\subseteq I_0=[0,n]\).
Write \(\ell_a=\beta_a-\alpha_a\).
The grouped description uses endpoints
\(t_jn/L,(L-t_j)n/L\), with multiplicity \(\nu_j\).

Set
\begin{equation}\label{eq:QandC}
\begin{split}
Q_n(z)&=\prod_{a=0}^{m}\prod_{r=\alpha_a}^{\beta_a}(z+r),\\
N_n&=\sum_{a=0}^{m}\ell_a,\qquad
d_n=\lcm(1,\ldots,n),\\
C_n&=m!\,d_n^m\prod_{a=0}^{m}\ell_a!.
\end{split}
\end{equation}
The degree of \(Q_n\) is \(N_n+M\).
\emph{The interval factorials are in the numerator of \(C_n\).}

Let \(\Gamma\) positively surround the poles \(0,-1,\ldots,-n\).
On a local logarithm branch define
\begin{equation}\label{eq:Rdef}
R_{h,n}(x)=\frac{C_n}{2\pi\ii}
\oint_\Gamma \frac{z^h x^{z+n}}{Q_n(z)}\dd z,\qquad 0\le h\le m.
\end{equation}
Write
\[
\frac{z^h}{Q_n(z)}
=\sum_{\ell=0}^{n}\sum_{r=1}^{\mu_\ell}
\frac{b_{h,\ell,r}}{(z+\ell)^r},
\quad
\mu_\ell=\#\{a:\ell\in I_a\}.
\]
Residues give
\begin{align}
R_{h,n}(x)&=\sum_{j=0}^{m}A_{h,j,n}(x)(\log x)^j,\label{eq:Rexpand}\\
A_{h,j,n}(x)&=\frac{C_n}{j!}
\sum_{\ell:\mu_\ell\ge j+1}
b_{h,\ell,j+1}x^{n-\ell}.\label{eq:Adef}
\end{align}

\subsection{Why all coefficients are integers}

\begin{lemma}[Integrality]\label{lem:integrality}
For all \(h,j\), one has \(A_{h,j,n}(x)\in\ZZ[x]\).
\end{lemma}
\begin{proof}
Absorb the interval factorials first:
\[
H_a(z)=\frac{\ell_a!}{\prod_{r=\alpha_a}^{\beta_a}(z+r)}.
\]
If \(\ell\in I_a\), expansion at \(z=-\ell+w\) gives
\begin{equation}\label{eq:polelocal}
wH_a(-\ell+w)=
(-1)^{\ell-\alpha_a}\binom{\ell_a}{\ell-\alpha_a}
\prod_{\substack{r\in I_a\\r\ne\ell}}
\left(1+\frac{w}{r-\ell}\right)^{-1}.
\end{equation}
The coefficient of \(w^u\) has denominator dividing \(d_n^u\):
every distance is a nonzero integer of absolute value at most \(n\).

If \(\ell\notin I_a\), use
\begin{equation}\label{eq:regularpartial}
H_a(z)=\sum_{r=\alpha_a}^{\beta_a}
\frac{(-1)^{r-\alpha_a}\binom{\ell_a}{r-\alpha_a}}{z+r}.
\end{equation}
Now the coefficient of \(w^u\) has denominator dividing \(d_n^{u+1}\).
The extra factor matters: a regular factor need not have an integral
constant term.

There are \(\mu_\ell\) pole factors and \(M-\mu_\ell\) regular factors.
The coefficient of \((z+\ell)^{-j-1}\) requires total local order
\(\mu_\ell-1-j\). The polynomial \((-\ell+w)^h\) has integer
coefficients and can only reduce the order required from the other
factors. Thus the necessary exponent of \(d_n\) is at most
\[
(\mu_\ell-1-j)+(M-\mu_\ell)=m-j.
\]
Multiplication by \(m!d_n^m/j!\) clears the denominators.
\end{proof}

The cancellation of \(\mu_\ell\) is why a variable pole multiplicity
does not require a larger universal power than \(d_n^m\).

\subsection{Supports and vanishing orders}

Nestedness implies
\[
\mu_\ell\ge j+1\quad\Longleftrightarrow\quad
\ell\in[\alpha_j,\beta_j].
\]
Hence \(A_{h,j,n}\) is divisible by \(x^{n-\beta_j}\), and the
remaining degree is at most \(\ell_j\).

\begin{lemma}[Vanishing at \(1\)]\label{lem:vanishing}
Every \(R_{h,n}\) has a zero of order at least \(N_n\) at \(x=1\).
\end{lemma}
\begin{proof}
Put \(w=\log x\). The coefficient of \(w^r\) in the integral of
\(z^h e^{zw}/Q_n(z)\) is zero whenever \(h+r\le\deg Q_n-2\), by the
residue at infinity. The first possible order is therefore
\[
\deg Q_n-h-1=N_n+M-h-1\ge N_n.
\]
Multiplication by \(x^n\) does not change this order.
Also \(\log x=(x-1)+O((x-1)^2)\).
\end{proof}

\subsection{A nonzero determinant, including its leading term}

\begin{proposition}[Exact determinant]\label{prop:det}
There exists \(c_n\ne0\) such that
\begin{equation}\label{eq:detidentity}
\det(A_{h,j,n}(x))_{0\le h,j\le m}
=c_n x^{\sum_j(n-\beta_j)}(x-1)^{N_n}.
\end{equation}
In particular \(A_n(\ii)\) is invertible.
\end{proposition}
\begin{proof}
The supports give the power of \(x\). After it is removed, the
determinant has degree at most \(N_n\).
Replace column zero by the sum of column \(j\) multiplied by
\((\log x)^j\), for \(0\le j\le m\).
Its determinant is unchanged; the new column consists of the
\(R_{h,n}(x)\). Lemma~\ref{lem:vanishing} supplies the factor
\((x-1)^{N_n}\).

To prove nonvanishing, take a block of repeated left endpoints
\[
\alpha_u=\cdots=\alpha_v=\alpha,\qquad r=v-u+1.
\]
The pole at \(-\alpha\) has multiplicity \(v+1\). Write
\[
\frac1{Q_n(-\alpha+w)}
=w^{-(v+1)}B_\alpha(w),\qquad B_\alpha(0)\ne0.
\]
For \(u\le j\le v\), the coefficient of \(x^{n-\alpha}\) in column
\(j\) is
\[
\frac{C_n}{j!}[w^{v-j}]
\bigl((-\alpha+w)^hB_\alpha(w)\bigr).
\]
The orders \(v-j\) are \(0,\ldots,r-1\), in reverse order.
Multiplication by \(B_\alpha\) gives an invertible triangular
change on these functionals.

In all blocks together, the leading-coefficient matrix reduces to
the divided-derivative evaluations of \(1,z,\ldots,z^m\) at distinct
points \(-\alpha\). This confluent Vandermonde matrix is invertible:
a polynomial in its kernel would have zeros of total multiplicity
\(M=m+1\), exceeding its degree. Thus the leading determinant is
nonzero, proving the identity.
\end{proof}

High vanishing order and a degree bound alone would also allow the
zero determinant. The leading-term argument rules out that possibility.
